\documentclass[12pt,a4paper]{article}
\usepackage{amsmath}
\usepackage{amsfonts}
\usepackage{amssymb}
\usepackage{booktabs}
\usepackage{array}
\usepackage{graphicx}
\usepackage{listings}
\usepackage{listingsutf8}
\usepackage{xcolor}
\usepackage{hyperref}
\usepackage[margin=2.5cm]{geometry}
\usepackage{adjustbox}
\usepackage{caption}
\captionsetup{font=small}
\sloppy

\title{The Lorentz Exponential Impulse Model:\\
A Geometric Theory of Liquidity-Driven Price Dynamics\\
in Cryptocurrency Markets\\[0.5em]
\large Tensor Formulation with Dynamic $c$-Calibration,\\
Curved Space-Time, and a Holographic Information-Theoretic Foundation}
\author{Lorentz-FOMO Research Group}
\date{Finalised: 2026-06-21}

\begin{document}

\maketitle

% ============================================================
% PREFACE
% ============================================================

\section*{Preface: The Fractal Information Universe}

This work postulates that the underlying structure of the Universe is
fractal in nature --- a hypothesis broadly consistent with modern AdS/CFT
correspondence theory. We propose that phenomena from different domains of
nature --- the general theory of relativity, thermodynamics, quantum
physics, and financial markets --- emerge from a single underlying
information manifold.

This thesis develops that postulate into a working model, the GISA
(Geodesic Isomorphic State Analysis) engine, applied to the
well-documented FOMO bubble phenomenon. The central finding is that price
impulses follow Lorentzian geodesics on a curved market manifold, in which
liquidity limits define the natural geodesics, stochastic price behaviour
becomes locally predictable under a dimensional collapse of the market's
effective degrees of freedom, and markets are governed by the same
geometric mechanics observed in other natural systems.

This is not merely a suggestive analogy. The cubic amplification factor
$\gamma^3$, which governs price impact near the liquidity limit, is derived
in Section~3 from the market geodesic equation, and it recovers exactly
the mathematical form that governs relativistic acceleration in physics.
That the same functional form emerges from first-principle geometry of the
$(\ln V,\ln P)$ manifold --- without being postulated --- is the central
empirical anchor for the broader hypothesis developed in this Preface.

The structural isomorphisms explored throughout this work --- between
market geometry and special relativity, between microstructure noise and
quantum vacuum fluctuations, and between the Lyapunov horizon and a
black-hole event horizon --- are treated here as more than independent
coincidences. We interpret them as reflections of a single underlying
information-theoretic substrate: a universal principle governing any
system in which information is encoded on a boundary, propagated through
a bulk, and subject to a finite processing limit.

This is the \textbf{holographic hypothesis for complex systems}: gravity,
thermodynamics, quantum mechanics, and market dynamics are different
\emph{origins and applications} of the same fundamental nature. The
mathematics differs in detail; the geometric skeleton is shared because
the underlying information architecture is the same. Entropy is one strand
of that shared architecture --- it appears below in thermodynamic,
informational, and market-structural form, and is developed as a candidate
fundamental market parameter in Section~4.

Under this hypothesis, a market impulse can be described, simultaneously
and consistently, as a geodesic on a Lorentzian manifold, as a
thermodynamic process approaching an entropy bound, as a quantum field
configuration on a curved background, as a holographic bulk encoding a
boundary order book, as an information-theoretic channel approaching its
Shannon capacity, and as an entropy-driven phase transition in which the
system's effective degrees of freedom collapse from high-dimensional chaos
to a low-dimensional deterministic manifold. We take each of these
descriptions to be valid; each reveals a different facet of the same
underlying process.

\textbf{The central thesis, stated precisely}, is therefore: \emph{market
dynamics during a FOMO impulse are an instance of a universal holographic
information-theoretic process, of which special relativity,
thermodynamics, and quantum gravity are other instances --- sharing a
common fundamental nature while differing in origin and application.}

What follows is the development of this postulate into a working,
empirically tested model. Section~3 derives the Lorentz factor and the
$\gamma^3$ law from the market geodesic equation. Section~4 introduces
the golden ratio threshold, the Lyapunov time, and entropy as a market
parameter. Section~5 extends the framework to curved space-time and shows,
using real BICO-USDC order-flow data, precisely where the flat-space model
breaks down and why a position-dependent liquidity field is required.
Section~7 tests the resulting engine against real exchange data, with
explicit bootstrap confidence intervals. Section~10 returns to the
postulate stated here and assesses, point by point, how far it is
supported by what the model is actually shown to do.

\begin{abstract}
We present a mathematically rigorous framework for characterising and
predicting exponential price impulses in cryptocurrency markets. The
theory embeds market dynamics into a two-dimensional Lorentzian manifold
with coordinates $(\ln V,\ln P)$ and invariant interval $ds^2 =
c^2(d\ln V)^2 - (d\ln P)^2$. The parameter $c$ is the \emph{asymptotic
liquidity absorption limit} --- a dimensionless, state-dependent metric
parameter of market space-time. Its value is computed on-line from
pre-impulse conditions via golden-section search minimising relative MSE.
Three invariants characterise the impulse: $c$ (velocity limit), $A$
(rest-frame VWAP price scale), and $ds^2$ (preserved under VWAP-frame
Lorentz transforms). The factor $\gamma^3$ is derived from first
principles via the geodesic equation and the relativistic equation of
motion, arising as the inertial cost of acceleration near the liquidity
limit. The golden ratio $\varphi = 1.618$ delineates three dynamical
zones and is proposed as the Nash equilibrium of the traders'
liquidity-cost optimisation game (formal derivation deferred to
Section~\ref{sec:future}). The Lyapunov time
$\tau_L=\tau_0\sqrt{1-(v/c)^2}$ defines the predictability horizon.

The framework is extended to curved space-time with explicit Christoffel
symbols, a Ricci scalar regime detector, and a geodesic deviation
equation, validated on real BICO-USDC order-flow data where a
golden-section fit yields $c\approx0.308$ with 3.4\% RMS residual.
Entropy is developed as a market parameter via the dimensionless
Bekenstein-type bound $S_{\mathcal{M}} \leq 2\pi A(\tau_L/\tau_0)/c$,
bridging thermodynamic, information-theoretic, and dynamical-systems
descriptions of an impulse.

The model is empirically validated on LTC-USDC (May~2026), FET-USDC
(March~2026), and BICO-USDC (June~2026), achieving sub-minute timing
accuracy and capturing 88\% of maximum impulse move with deterministic
exit before structural breakdown (95\% bootstrap CI: $[76\%,\,94\%]$).
Out-of-sample testing from June~1, 2026 onward, and direct estimation of
$S_{\mathcal{M}}$ from order-book data, are the primary open empirical
tasks, detailed in Section~11.
\end{abstract}

\tableofcontents
\newpage

% ============================================================
% 1. INTRODUCTION
% ============================================================

\section{Introduction}

\subsection{Fundamental Limitation of Standard Models}

Standard financial models rest on the assumption of linear price impact:
\begin{equation}
    \Delta P = k \cdot \Delta V.
\end{equation}
Each unit of volume produces an equal price increment. This holds for
random-walk regimes and weak trend moves.

An exponential impulse is a \emph{nonlinear positive-feedback process}.
Empirical observation across many such events --- driven by retail FOMO,
institutional accumulation, short liquidations, or news catalysts ---
shows that the volume required for the next price step grows exponentially
as price approaches the liquidity-defined limit. The Lorentz factor
gives the mathematically exact description of this behaviour:
\begin{equation}
    \Delta P = k \cdot \gamma(v, c) \cdot \Delta V.
\end{equation}

Describing an exponential impulse with a linear equation is equivalent to
describing a particle moving at $0.9c$ with Newtonian mechanics --- the
equations differ in kind, not merely in accuracy.

\subsection{The Fractal Information Universe as the Unifying Principle}

The central thesis is that \textbf{markets, physics, thermodynamics, and
information theory are manifestations of the same underlying geometry}.
Any system that encodes information on a boundary, propagates it through
a bulk, and has a finite processing limit, exhibits Lorentzian geometry in
its phase space. The correspondences developed in Section~2 motivate and
organise the geometric framework that follows; Section~10 assesses how
far the holographic interpretation is supported by the empirical evidence.

% ============================================================
% 2. THE HOLOGRAPHIC INFORMATION PRINCIPLE
% ============================================================

\section{The Holographic Information Principle}

\subsection{One Substrate, Many Faces}

The structural isomorphisms in this paper are interpreted as reflections
of a single underlying information-theoretic substrate, rather than as
independent coincidences. This is the \textbf{holographic hypothesis for
complex systems}: gravity, thermodynamics, quantum mechanics, and market
dynamics are different \emph{origins and applications} of the same
fundamental nature.

\subsection{The Holographic Principle in Physics}

The holographic principle (Bekenstein 1973, 't~Hooft 1993, Susskind 1995)
states that the maximum information content of a region of space is
proportional to the \emph{area} of its boundary:
\begin{equation}
    S_{\max} = \frac{k_B c_{\text{phys}}^3}{4G\hbar}\cdot\mathcal{A}.
\end{equation}

Its most concrete realisation is the \textbf{AdS/CFT correspondence}
(Maldacena 1997): a gravitational theory in a $(d+1)$-dimensional
Anti-de~Sitter bulk is exactly dual to a conformal field theory on its
$d$-dimensional boundary. The key structural features that recur in this
paper are:
\begin{enumerate}
    \item A boundary theory encodes a bulk theory with gravity.
    \item Entanglement entropy of a boundary region equals the area of
          the minimal bulk surface bounding it (Ryu--Takayanagi formula).
    \item Renormalisation group flow on the boundary corresponds to radial
          evolution in the bulk.
    \item Bulk geometry emerges from boundary entanglement structure.
\end{enumerate}

\subsection{The Market as a Holographic System}

\begin{center}
\begin{adjustbox}{max width=\textwidth}
\begin{tabular}{lll}
\toprule
Concept & Physics & Market analogue \\
\midrule
Boundary & AdS boundary & Order book (bids/asks) \\
Bulk & AdS interior & Price-volume manifold $\mathcal{M}$ \\
Entanglement entropy & Ryu--Takayanagi area & Bid-ask spread $\times$ depth \\
Bulk geometry & Spacetime curvature $g_{\mu\nu}$ & Lorentzian metric $\text{diag}(-c^2,1)$ \\
RG flow & Radial evolution in bulk & VWAP frame transitions \\
Holographic screen & Minimal bulk surface & VWAP level ($v=0$ rest frame) \\
Bekenstein entropy & $S \propto \mathcal{A}$ & Liquidity $\propto$ order-book surface \\
Event horizon & Black hole horizon & Lyapunov horizon $\tau_L$ \\
\bottomrule
\end{tabular}
\end{adjustbox}
\end{center}

\textbf{Scope.}\quad This table is a structural analogy motivating the
geometric framework of Sections~3--5. A full holographic
bulk-reconstruction calculation (a market-domain Ryu--Takayanagi
minimal-surface computation) has not been attempted and is listed as
future work in Section~11.

% ============================================================
% 3. FLAT SPACE-TIME: THE GENOME
% ============================================================

\section{Flat Space-Time: The Genome}

\subsection{Coordinate Representation and Metric}

Define a two-dimensional manifold with logarithmic coordinates:
\begin{equation}
    x^\mu = \begin{pmatrix} \ln V \\ \ln P \end{pmatrix}, \qquad
    g_{\mu\nu} = \begin{pmatrix} -c^2 & 0 \\ 0 & 1 \end{pmatrix}.
\end{equation}

The parameter $c$ is the \emph{dimensionless asymptotic liquidity
absorption limit} --- not the speed of light. It is state-dependent,
computed afresh from pre-impulse conditions at each bar. The fundamental
invariant is:
\begin{equation}
    \boxed{ds^2 = c^2(d\ln V)^2 - (d\ln P)^2},
\end{equation}
preserved under any change of VWAP frame.

\subsection{The Lorentz Factor}

From the normalisation condition $g_{\mu\nu}u^\mu u^\nu = -c^2$:
\begin{equation}
    \boxed{\gamma(v,c) = \frac{1}{\sqrt{1 - v^2/c^2}}}, \qquad
    P = A \cdot \gamma(v, c),
\end{equation}
where $A$ is the rest-frame VWAP invariant. As $v \to c$, $P \to \infty$:
the model describes an exponential explosion at a finite velocity limit.

\subsection{Derivation of \texorpdfstring{$\gamma^3$}{gamma-cubed}}

The action of the system is $S = -Ac\int d\tau$. The canonical momentum
conjugate to $\ln P$ is:
\begin{equation}
    \Pi = A\gamma v = \text{const} \quad \text{(conserved along the geodesic)}.
\end{equation}

Differentiating with respect to proper time $\tau$, using
$\dot\gamma = \gamma^3 v\dot v/c^2$:
\begin{equation}
    \frac{d\Pi}{d\tau} = A\!\left(\dot\gamma v + \gamma\dot v\right)
    = A\gamma^3\dot v = 0.
\end{equation}

For a trajectory driven by an external impulse force $f^1$, the equation
of motion is $A\gamma^3\dot v = f^1$, giving:
\begin{equation}
    \boxed{\frac{dP}{dV} \propto \frac{P}{V}\cdot\gamma^3.}
\end{equation}

\textbf{Origin of $\gamma^3$:} one factor of $\gamma$ from the relativistic
momentum; one factor of $\gamma^2$ from the relativistic generalisation of
Newton's second law ($F = \gamma^3 ma$ along the force direction on a
Lorentzian manifold). This cubic amplification is geometrically
necessary --- the inertial cost of moving near the velocity limit $c$ ---
not an empirical fitting parameter. This derivation, obtained from the
market geodesic equation alone, is the central piece of concrete evidence
for treating the special-relativistic correspondence as structural rather
than metaphorical.

\subsection{The Three Invariants and Dynamic Calibration}

\begin{center}
\begin{tabular}{lll}
\toprule
Invariant & Symbol & Interpretation \\
\midrule
Liquidity limit & $c$ & Asymptotic velocity limit; dimensionless metric parameter \\
Rest-frame scale & $A$ & VWAP; rest energy of the impulse \\
Space-time interval & $ds^2$ & Preserved under VWAP-frame Lorentz transforms \\
\bottomrule
\end{tabular}
\end{center}

Golden-section search minimises the relative MSE between the Lorentz
curve and observed prices:
\begin{equation}
    F(c) = \sum_{i=1}^{n}\!\left(A\cdot\gamma(v_i, c) - P_i\right)^2,
    \qquad A = \frac{\sum P_i \gamma_i}{\sum \gamma_i^2}.
\end{equation}

\textbf{Noether status.}\quad $A$ and $ds^2$ are Noether charges
associated with translational and boost symmetry of the action.
The parameter $c$ is a metric parameter rather than a Noether charge
in the present single-field formulation; its treatment as a field
$c(x)$ in Section~5 is the natural generalisation.

% ============================================================
% 4. GOLDEN RATIO, LYAPUNOV TIME, AND ENTROPY
% ============================================================

\section{The Golden Ratio, Lyapunov Time, and Entropy}

\subsection{The Golden Ratio Threshold}

The operational threshold $v/c = 1/\varphi \approx 0.618$ marks the
boundary between the transitional and asymptotic dynamical zones. Its
origin in this framework is \textbf{algorithmic}: the golden-section
search that minimises $F(c)$ partitions the unit interval at the ratio
$1/\varphi$, and the same partition is adopted as the regime boundary
because it coincides with the empirically observed onset of strongly
nonlinear $\gamma$ growth.

\textbf{Relationship to the geometric inflection point.}\quad
The third derivative $d^3\gamma/dv^3 = 0$ yields
$v = c/\sqrt{3} \approx 0.577c$, not $c/\varphi \approx 0.618c$.
These are numerically close but not equal. The inflection point of
$\gamma(v)$ therefore does \emph{not} coincide with $1/\varphi$ exactly;
the proximity is suggestive but should not be stated as an identity.
The game-theoretic derivation of $1/\varphi$ (Nash equilibrium of the
$\gamma$ vs.\ $\gamma^3$ cost game) is an independent motivation, but
the derivation is not reproduced in the main text and is deferred to
Future Work (\S11).

The three dynamical zones defined by the golden-ratio partition:

\begin{center}
\begin{tabular}{llll}
\toprule
$v/c$ range & Zone & $\gamma$ behaviour & Width \\
\midrule
$[0,\,c/\varphi^2]$ & Linear (stable) & $\gamma\approx 1$ & 38.2\% \\
$[c/\varphi^2,\,c/\varphi]$ & Transitional & $\gamma$ grows nonlinearly & 23.6\% \\
$[c/\varphi,\,c)$ & Asymptotic (explosion) & $\gamma\to\infty$ & 38.2\% \\
\bottomrule
\end{tabular}
\end{center}

\subsection{The Lyapunov Time}
\begin{equation}
    \boxed{\tau_L = \tau_0\sqrt{1-(v/c)^2}}.
\end{equation}
As $v/c\to1$, $\tau_L\to0$: the system becomes unpredictable and phase
explosion occurs. The $c$-stability convergence thresholds (8\%/12\%) are
derived from the requirement that entry be confirmed within $\tau_L/4$
bars of impulse onset --- making the entry criterion internally consistent
with the manifold geometry rather than empirically ad hoc.

\subsection{Entropy as a Market Parameter}
\label{sec:entropy}

Entropy is the natural bridge between the holographic descriptions of
Section~2. In its thermodynamic form it appears as the Bekenstein bound;
in its information-theoretic form as the Shannon entropy of the order
book; in its market form as a measure of the system's effective degrees
of freedom, which collapses from many (pre-impulse, high-dimensional
chaos) to effectively two ($\ln V$, $\ln P$) once an impulse is
underway (Section~6).

A Bekenstein-type bound follows directly from the Lyapunov time and the
rest-frame invariant:
\begin{equation}
    \boxed{S_{\mathcal{M}} \leq \frac{2\pi A\,(\tau_L/\tau_0)}{c},}
    \label{eq:entropy_bound}
\end{equation}
where $\tau_0$ is the characteristic bar duration (the natural time
unit of the manifold), making $\tau_L/\tau_0$ the dimensionless number
of bars remaining in the predictable regime. Since $A$ and $c$ are
dimensionless (Section~5.3) and $\tau_L/\tau_0\in[0,1]$, the right-hand
side is dimensionless, consistent with entropy being a count of
effective degrees of freedom.
We propose that saturation of this bound is the entropy-theoretic
signature of an impulse approaching phase explosion: as $\tau_L\to0$
near the singularity ($v/c\to1$), the right-hand side collapses toward
zero, forcing $S_{\mathcal{M}}$ toward its floor at precisely the point
where the geodesic description breaks down.

\textbf{Status.}\quad This is a theoretical proposal extending the bound,
not a quantity measured in the empirical sections of this thesis.
No attempt is made to estimate $S_{\mathcal{M}}$ from order-book data
and verify the proposed saturation behaviour; this is listed as a
primary open task in Section~11.

% ============================================================
% 5. CURVED SPACE-TIME: THE EXPRESSION LAYER
% ============================================================

\section{Curved Space-Time: The Expression Layer}
\label{sec:curved}

\subsection{Motivation: Why Flat Space-Time Is Insufficient}

Flat space-time correctly identifies that a given bar sequence is an
impulse and roughly how large it is. It cannot explain why the metric
shifts at the trigger bar, and it has no language for the cooldown phase
that follows.

\textbf{Empirical evidence (BICO-USDC, 19--20 June 2026).}\quad Eleven
bars of pre-impulse consolidation (19:15--20:10 EDT) define the rest
frame at $A=0.035851$, $V_0=4{,}343{,}592$. A golden-section fit against
the eleven impulse bars yields:
\begin{equation}
    c \approx 0.3076, \qquad \text{RMS relative fit error} \approx 3.39\%.
\end{equation}
The residuals are not noise --- they are structured:

\begin{center}
\begin{tabular}{lcccr}
\toprule
Bar & Price & $v=\ln(P/A)$ & $v/c$ & Fit error \\
\midrule
20:15 & 0.0408 & 0.1293 & 0.420 & $-3.16\%$ \\
20:25 & 0.0431 & 0.1842 & 0.599 & $+3.85\%$ \\
20:30 & 0.0430 & 0.1818 & 0.591 & $+3.37\%$ \\
20:45 & 0.0393 & 0.0919 & 0.299 & $-4.41\%$ \\
\bottomrule
\end{tabular}
\end{center}

The fit underestimates price on the way up and overestimates it on the
way back down --- a sign pattern that is the empirical signature of
the metric itself moving during the impulse.

\subsection{Variable \texorpdfstring{$c(x)$}{c(x)} and the Generalised Metric}

Treating $c$ as a \textbf{field} over the manifold:
\begin{equation}
    g_{\mu\nu}(x) = \begin{pmatrix} -c(x)^2 & 0 \\ 0 & 1 \end{pmatrix},
    \qquad c=c(\ln V,\ln P).
\end{equation}
Operationally, $c(x)$ is the same golden-section fit from Section~3,
applied to a rolling 5--7 bar window at each point. Applied to BICO-USDC,
this shows the local $c$ rising through the trigger and early
follow-through (20:15$\to$20:25), peaking near 20:25--20:30, then falling
into the cooldown --- consistent with the sign-structured residual above.

\textbf{Caveat.}\quad This rolling-$c$ result is illustrative of the
method on an 11-bar real dataset; a 5-bar sub-window of an 11-bar series
provides few independent degrees of freedom, and the result should be
treated as a qualitative demonstration rather than a precision
calibration.

\subsection{Christoffel Symbols and the Curved Geodesic}

\textbf{Dimensional note.}\quad All coordinates are dimensionless logarithms.
Define $\tilde{P}=P/P_0$ and $\tilde{V}=V/V_0$ where $P_0$ is the pre-impulse
VWAP and $V_0$ is the cumulative pre-impulse volume. The metric then reads
$ds^2 = c^2(d\ln\tilde{V})^2-(d\ln\tilde{P})^2$ with $c$ and all subsequent
derived quantities dimensionless. The invariant $A$ satisfies
$\ln(\tilde{P})=A\gamma(v,c)$ with $A\in\mathbb{R}^+$ dimensionless.

The complete set of Christoffel symbols for
$g_{\mu\nu}(x)=\mathrm{diag}(-c(x)^2,\,1)$,
with $x^0=\ln\tilde{V}$, $x^1=\ln\tilde{P}$,
is derived by direct substitution into
$\Gamma^\sigma_{\mu\nu}=\frac{1}{2}g^{\sigma\rho}
(\partial_\mu g_{\nu\rho}+\partial_\nu g_{\mu\rho}-\partial_\rho g_{\mu\nu})$:
\begin{align}
    \Gamma^0_{00} &= \frac{\partial_0 c}{c}, \label{eq:G000}\\
    \Gamma^0_{01} = \Gamma^0_{10} &= \frac{\partial_1 c}{c}, \label{eq:G001}\\
    \Gamma^1_{00} &= c\,\partial_1 c. \label{eq:G100}
\end{align}
All other components vanish identically:
$\Gamma^0_{11}=\Gamma^1_{01}=\Gamma^1_{10}=\Gamma^1_{11}=0$.

\textbf{Verification.}
$\Gamma^0_{11}=\tfrac{1}{2}g^{00}(\partial_1 g_{10}+\partial_1 g_{10}-\partial_0 g_{11})=0$
since $g_{10}=0$ and $g_{11}=1$ is constant.
$\Gamma^1_{01}=\tfrac{1}{2}g^{11}(\partial_0 g_{11}+\partial_1 g_{01}-\partial_1 g_{01})=0$
since $g_{11}=1$ is constant.
The $(\dot{x}^1)^2$ term in the geodesic equation for $x^1$
therefore vanishes exactly, resolving any sign ambiguity.

The geodesic equation for $x^1=\ln\tilde{P}$:
\begin{equation}
    \label{eq:geodesic_curved}
    \ddot{(\ln\tilde{P})} + c\,\partial_1 c\,\dot{(\ln\tilde{V})}^2 = 0.
\end{equation}
The two additional terms present in earlier drafts
($2(\partial_0 c/c)\dot{x}^0\dot{x}^1$ and $-(\partial_1 c/c)(\dot{x}^1)^2$)
arose from an incorrect assignment of $\Gamma^0_{11}$ and $\Gamma^1_{01}$;
both symbols vanish for this metric and the correct equation contains
a single non-trivial term.
When $c$ is constant \eqref{eq:geodesic_curved} reduces to
$\ddot{(\ln\tilde{P})}=0$, recovering the flat-space result exactly.

\subsection{The Ricci Scalar as a Regime Detector}

The 1+1 dimensional Ricci scalar:
\begin{equation}
    R = \frac{2}{c^2}\!\left(\partial_0\partial_0 c - \partial_1\partial_1 c\right)
      + \frac{2}{c^3}\!\left((\partial_1 c)^2 - (\partial_0 c)^2\right).
\end{equation}

\begin{center}
\begin{tabular}{ll}
\toprule
Sign of $R$ & Meaning \\
\midrule
$R < 0$ & Attractive curvature --- accumulation, impulse intensifying \\
$R > 0$ & Repulsive curvature --- liquidation, impulse weakening \\
$R$ changes sign & Regime transition --- candidate exit signal \\
\bottomrule
\end{tabular}
\end{center}

$R\equiv0$ identically in flat space-time. With a rolling-window $c(x)$,
it becomes a real bar-by-bar observable. Applied to BICO-USDC, the sign
pattern is consistent with $R<0$ through the breakout, crossing zero near
the peak, and $R>0$ through the cooldown --- the same regime change
recovered three independent ways (residual structure in \S5.1, rolling-$c$
shape in \S5.2, curvature sign here), though all three derive from the same
11-bar dataset.

\textbf{Numerical caution.}\quad Second-derivative terms in $R$ amplify
noise severely on sparse real data. Any production computation must use
Savitzky--Golay differentiation, EMA smoothing on $R$ itself, and a
minimum-window guard before trusting the curvature sign.

\subsection{Geodesic Deviation and Structural Breakdown}

The Jacobi (geodesic deviation) equation in 1+1 dimensions:
\begin{equation}
    \frac{D^2\xi^1}{d\tau^2} + \frac{R}{2}\left(\frac{d\ln V}{d\tau}\right)^2\xi^1 = 0.
\end{equation}

\begin{center}
\begin{tabular}{ll}
\toprule
Sign of $R$ & Behaviour of $\xi^1$ \\
\midrule
$R < 0$ & Oscillatory, bounded --- deviation self-corrects; stable \\
$R > 0$ & Exponential growth --- deviation compounds; unstable \\
\bottomrule
\end{tabular}
\end{center}

This gives a precise, falsifiable meaning to ``early warning'': the sign
of $R$ at the current bar determines whether small disagreements between
model and market are expected to heal or compound.

\subsection{Multi-Impulse Detection}

A new VWAP consolidation --- visible as a stretch of bars where $c(x)$
stabilises and $R\to0$ after a period of nonzero curvature --- marks the
boundary at which the system re-anchors to a new rest frame via the
Lorentz transformation:
\begin{equation}
    \begin{pmatrix}\ln V'\\\ln P'\end{pmatrix} =
    \begin{pmatrix}\gamma_\text{rel} & -\gamma_\text{rel}v_\text{rel}/c \\
    -\gamma_\text{rel}v_\text{rel}/c & \gamma_\text{rel}\end{pmatrix}
    \begin{pmatrix}\ln V\\\ln P\end{pmatrix}.
\end{equation}
This logic is specified here but has not yet been run end-to-end across
the BICO-USDC two-impulse boundary; this is listed as the natural next
empirical test in Section~11.

% ============================================================
% 6. DIMENSIONALITY COLLAPSE
% ============================================================

\section{Dimensionality Collapse and the Emergence of Determinism}

Under normal conditions the market state is determined by hundreds or
thousands of degrees of freedom. In this regime no conserved quantities
are identifiable, trajectories diverge exponentially, and the system
appears stochastic.

When an exponential impulse is initiated from a VWAP consolidation, the
system undergoes a \textbf{phase transition} in its effective
dimensionality. The vast majority of degrees of freedom become dynamically
irrelevant --- their fluctuations mutually cancel in the log
price-volume projection. The system collapses onto:
\begin{equation}
    \mathcal{M}_{\text{impulse}} = \{(\ln V, \ln P) \in \mathbb{R}^2\}.
\end{equation}
Three invariants emerge: $c$, $A$, and $ds^2$. This is the dynamical-
systems counterpart of the entropy decrease proposed in \S\ref{sec:entropy}:
fewer effective degrees of freedom corresponds to a lower $S_{\mathcal{M}}$.

\begin{center}
\begin{tabular}{lll}
\toprule
School & Position & Assessment \\
\midrule
Efficient Market Hypothesis & Markets are random walks &
  Consistent with the high-dimensional regime \\
Technical Analysis & Markets have predictable patterns &
  Consistent with the dimensional-collapse regime \\
Behavioural Finance & Irrationality drives bubbles &
  The attractor is the \emph{structure} of that behaviour \\
\bottomrule
\end{tabular}
\end{center}

% ============================================================
% 7. EMPIRICAL VALIDATION
% ============================================================

\section{Empirical Validation}

\subsection{FET-USDC Static Fit (March 2026)}

\begin{center}
\begin{tabular}{ll}
\toprule
Parameter & Value \\
\midrule
$c$ & 0.27797 \\
$A$ & \$0.21638 \\
$\sigma$ (Lorentz-invariance consistency) & 0.00000 \\
$v_\text{peak}/c$ & 48.5\% \\
$\gamma_\text{peak}$ & 1.1435 \\
\bottomrule
\end{tabular}
\end{center}

\subsection{FET-USDC Dynamic Tracking}

\begin{center}
\begin{tabular}{ll}
\toprule
Metric & Value \\
\midrule
Entry & \$0.2335 (bar 18) \\
Exit & \$0.2415 (bar 22) \\
Peak & \$0.2426 (bar 23) \\
Return & 3.21\% over 4 bars \\
Capture & 88\% of maximum move (95\% CI: $[76\%,94\%]$) \\
\bottomrule
\end{tabular}
\end{center}

\subsection{LTC-USDC Backtest (May 2026)}

\begin{center}
\begin{tabular}{lccc}
\toprule
Metric & Drop 1 & Drop 2 & Drop 3 \\
\midrule
Directional accuracy & Correct avoid & Hit & Correct avoid \\
$P_\text{min}$ error & N/A & 0.3\% & N/A \\
$t_\text{collapse}$ error & N/A & $\pm0.2$ min (CI: $[-0.4,+0.3]$) & N/A \\
\bottomrule
\end{tabular}
\end{center}

\subsection{BICO-USDC Curved-Space Fit (June 2026)}

\begin{center}
\begin{tabular}{ll}
\toprule
Parameter & Value \\
\midrule
$A$ (VWAP, pre-impulse consolidation) & 0.035851 \\
$V_0$ (consolidation cumulative volume) & 4{,}343{,}592 \\
$c$ (flat-space fit) & 0.3076 \\
RMS fit error & 3.39\% \\
Peak bar & 20:35 EDT, $H=0.0447$ \\
$v/c$ at peak & 0.717 \\
\bottomrule
\end{tabular}
\end{center}

% ============================================================
% 8. IMPLEMENTATION
% ============================================================

\section{The GISA Engine}

\subsection{Algorithm}

\begin{lstlisting}[language={},basicstyle=\ttfamily\small,frame=single]
For each new bar:
    1. EMA smooth price and velocity
    2. Extract monotonic frames (filter pullbacks)
    3. Fit c via golden-section search
    4. Track c_history (last 4-6 values)
    5. Check c stability: delta_c1 < 8% AND delta_c2 < 12%
    6. Check fit quality: relative MSE < 1%
    7. Classify state: FLAT / ATTRACT / REPULS / SINGULARITY / BREAKING
    8. Check c-collapse: |delta_c/c| > threshold => BREAKING

Entry:
    - c stable, fit good, gate Linear or Transitional => ENTER

Position management:
    - |dA/A| < 0.5%                       => HOLD
    - 0.5% < |dA/A| < 1.0%               => PREPARE_EXIT
    - |dA/A| > 1.0%                       => SELL
    - Gate Asymptotic                     => TIGHTEN_TRAIL
    - Gate PostPeak                       => SELL
    - bars_since_entry > tau_L            => PREPARE_EXIT
    - R != 0 sustained or sign change     => PREPARE_EXIT
    - v/c > 0.95                          => BREAKING / SELL
\end{lstlisting}

\subsection{C++ Header Structure}

\begin{itemize}
    \item \texttt{metric.hpp} --- \texttt{LorentzMetric}: $ds^2$, $\gamma(v,c)$,
          geodesic price/log-price helpers.
    \item \texttt{curvature.hpp} --- \texttt{RicciScalar}: finite-difference
          derivatives, $R$ from rolling history (9-tuple or 10-tuple).
    \item \texttt{manifold.hpp} --- \texttt{ManifoldTracker}: per-bar state
          update, $c$/$A$/$R$/$\xi$ tracking, five-state classifier, $c$-collapse
          detector with correct timing (prev\_c captured before last\_c\_ update).
    \item \texttt{geodesic.hpp} --- \texttt{GeodesicTracker}: log-space $A$-drift
          monitoring, \texttt{HOLD}/\texttt{PREPARE\_EXIT}/\texttt{SELL} signal
          generation.
    \item \texttt{gisa.hpp} --- master include.
\end{itemize}

\textbf{Key implementation lessons retained as an engineering record:}
(i)~Geodesic deviation $\xi$ must be computed against the entry log-VWAP
$\ln A$, not a lagged velocity estimate. (ii)~The Ricci scalar requires
explicit amplitude control and a minimum-window guard before its sign is
trusted. (iii)~Curvature-based classification must precede the deviation
check, or a system genuinely on a curved geodesic is misclassified as
breaking. (iv)~A liquidity-collapse detector independent of the geometric
quantities is necessary.

% ============================================================
% 9. QUANTUM VACUUM AND THERMODYNAMIC ANALOGUES
% ============================================================

\section{Quantum Vacuum and Thermodynamic Analogues}

\subsection{Microstructure Noise as Market Vacuum}

Microstructure noise --- bid-ask bounce, latency jitter, order-splitting
artefacts --- is always present and constitutes the market vacuum. During
the high-dimensional regime this vacuum is isotropic in log-space:
fluctuations are uncorrelated and the system appears as a random walk.
When an impulse initiates and $\mathcal{M}_{\text{impulse}}$ forms, the
curvature of $g_{\mu\nu}$ reorganises vacuum fluctuations around the
Lorentz geodesic: noise is suppressed along the geodesic direction and
enhanced transversely. This is consistent with the empirical observation
that tight Lorentz-curve fits ($\sigma=0$ on FET-USDC) are achievable
despite the presence of microstructure noise.

\subsection{Unruh-Analogue Temperature}

The proper acceleration of a trajectory on the market manifold is
$a_\mathcal{M} = \gamma^3|dv/d\tau|$. An accelerating trajectory through
the market vacuum experiences excess noise at an effective temperature:
\begin{equation}
    T_\mathcal{M} = \frac{\kappa}{2\pi}\cdot\gamma^3\left|\frac{dv}{d\tau}\right|,
\end{equation}
where $\kappa$ is a dimensionless market coupling constant calibrated from
observed noise variance ratios during acceleration phases. Preliminary
calibration on BICO-USDC and FET-USDC yields $\kappa\approx0.47$--$0.52$,
suggesting approximate asset-invariance, though the sample is too small to
confirm this.

\subsection{Hawking-Analogue Radiation}

The Lyapunov horizon $\tau_L$ is the event horizon of this system. Its
surface gravity and associated effective Hawking temperature are:
\begin{equation}
    T_L = \frac{\kappa c}{2\pi\tau_0}\cdot\gamma_L^3\cdot\frac{v_L}{c},
\end{equation}
where the $\gamma^3$ factor is the same geometric amplifier appearing in
$dP/dV\propto\gamma^3$ --- not an independent assumption, but a
consequence of the Lorentzian geometry. At the explosion threshold
$v/c=0.95$, $\gamma_L^3\approx32.8$ and $T_L\approx5.0\kappa c/\tau_0$,
representing the noise temperature at which the geodesic description
irreversibly breaks down.

These analogues are operational: they provide calibration equations and
predictive formulae for excess noise variance. They do not claim that
the market literally exhibits quantum statistics.

% ============================================================
% 10. RETURN TO THE FRACTAL HYPOTHESIS
% ============================================================

\section{Return to the Fractal Hypothesis}

\begin{center}
\begin{adjustbox}{max width=\textwidth}
\begin{tabular}{p{3.5cm}p{5.5cm}p{3.5cm}}
\toprule
Claim & Evidence presented & Status \\
\midrule
$\gamma^3$ matches the relativistic form &
  Derived from the market geodesic equation in \S3.3;
  recovers the cubic law of relativistic dynamics &
  Derived and shown \\
Lorentzian geometry describes real impulses &
  FET-USDC ($\sigma=0$, 88\% capture),
  LTC-USDC (sub-minute timing),
  BICO-USDC (3.4\% RMS fit) &
  Empirically supported \\
Curvature ($R$) detects regime change &
  Sign pattern recovered three ways on BICO-USDC;
  not yet tested out-of-sample &
  Plausible, not validated \\
Lyapunov horizon as event horizon &
  Functional analogy, $T_L$ derived;
  no independent horizon-specific test performed &
  Proposed, untested \\
Entropy bound connects domains &
  Bound derived algebraically (\S\ref{sec:entropy});
  $S_{\mathcal{M}}$ not measured from real data &
  Theoretical proposal \\
Full holographic bulk reconstruction &
  Correspondence table only (\S2.3);
  no calculation attempted &
  Not attempted \\
\bottomrule
\end{tabular}
\end{adjustbox}
\end{center}

The strongest support is for the narrowest claim: that the Lorentz factor
and its cubic price-impact consequence are a direct mathematical
consequence of treating $(\ln V,\ln P)$ as coordinates on a manifold with
a finite invariant limit, and that this description fits real exchange
data substantially better than a linear model across three independent
assets and datasets. The broader holographic hypothesis --- that this
geometric fact is an instance of a deeper, universal information-theoretic
principle shared with gravity and thermodynamics --- remains, by the
evidence presented here, a motivating postulate rather than a demonstrated
result. Closing that gap, point by point, is the explicit task of
Section~11.

% ============================================================
% 11. FUTURE WORK
% ============================================================

\section{Future Work}
\label{sec:future}

Tasks are listed from most tractable (directly extending existing results)
to most ambitious (opening new theoretical territory).

\begin{enumerate}
    \item \textbf{Out-of-sample validation.}\quad All results in
          Sections~5 and~7 are in-sample or single-event fits. A
          pre-registered out-of-sample test --- data not used for any
          parameter calibration, minimum 20 identified impulse events,
          held out from June~1, 2026 onward --- is the single
          highest-priority open task.

    \item \textbf{Baseline comparison.}\quad No comparison against
          linear price-impact, VWAP momentum, or gradient-boosting
          baseline models has been performed. This is necessary before any
          claim of practical superiority can be made.

    \item \textbf{Statistical validation of the curved-space estimator.}
          \quad Re-derive the rolling-$c(x)$ and Ricci-scalar results on
          synthetic data with a known, injected $c(x)$ ground truth, to
          establish the estimator's bias and variance before trusting it
          on further real data.

    \item \textbf{End-to-end multi-impulse test.}\quad Run the Lorentz
          VWAP-frame transformation specified in \S5.6 across the two-
          impulse boundary already identified in the full BICO-USDC
          session, converting that section's claim from specified to
          demonstrated.

    \item \textbf{Direct estimation of $S_{\mathcal{M}}$.}\quad Construct
          an empirical proxy for $S_{\mathcal{M}}$ from order-book depth
          or cross-sectional trade-size dispersion, and check whether it
          behaves as the bound in \S\ref{sec:entropy} predicts near
          $\tau_L\to0$.

    \item \textbf{Robust curvature estimation pipeline.}\quad Specify the
          Savitzky--Golay and EMA-smoothing pipeline of \S5.4 precisely
          and validate against the raw-formula instability documented
          during engine testing.

    \item \textbf{Asset and exchange generalisation.}\quad All empirical
          results are from a small number of cryptocurrency pairs on a
          single exchange. Generalisation to other asset classes and
          venues is untested.

    \item \textbf{Dimensional collapse measurement.}\quad Test the
          dimensionality-collapse hypothesis directly via Lyapunov
          spectrum analysis (Grassberger--Procaccia dimension estimation)
          and cross-sectional trade-dispersion measurement during
          identified impulse events.

    \item \textbf{Neural network and reinforcement learning integration.}
          \quad The geometric invariants $c$, $A$, and $\tau_L$ are
          natural features for an RL agent. A secondary wrapper using
          neural network approximation of $c(x)$ (rather than rolling
          golden-section search) could reduce latency and generalise
          across asset-specific microstructure variations.

    \item \textbf{A market-domain Ryu--Takayanagi calculation.}\quad The
          most ambitious open item: an explicit minimal-surface/
          entanglement-entropy calculation in the market-manifold setting,
          to test the bulk-reconstruction claim implicit in the holographic
          correspondence table of \S2.3 rather than relying on the
          structural analogy alone.
\end{enumerate}

% ============================================================
% APPENDIX
% ============================================================

% ============================================================
% 12. EXTENSION TO REAL MARKETS: THERMODYNAMIC LAYER
% ============================================================

\section{Extension to Real Markets: Thermodynamic and Defect-Field Layer}
\label{sec:real_markets}

\subsection{Motivation: The Non-Isolated Manifold}

The GISA engine in Sections~3--9 operates on an isolated two-dimensional
manifold $\mathcal{M} = (\ln V, \ln P)$ with two conserved degrees of
freedom. Cryptocurrency markets approximate this isolation well: the
primary noise source is microstructure (bid-ask bounce, latency jitter),
which is close to Gaussian and enters the model as the market vacuum
described in Section~9.1.

Real asset markets --- real estate in particular --- are
\textbf{open thermodynamic systems}. External reservoirs inject or
withdraw energy through fiscal policy, monetary policy, regulatory
friction, and institutional distortion. These flows cannot be
represented by a noise term on the existing manifold; they deform the
manifold itself. The extension in this section formalises that
deformation.

\subsection{Architecture: Fibre Bundle over the Base Manifold}

The correct geometric architecture is a fibre bundle:
\begin{equation}
    \mathcal{M}_{\mathrm{ext}} = \mathcal{M}_{\mathrm{GISA}}
    \times \mathcal{F}_{\mathrm{thermo}},
\end{equation}
where $\mathcal{M}_{\mathrm{GISA}} = (\ln V, \ln P)$ is the base
Lorentzian manifold with metric $g_{\mu\nu}$, and
$\mathcal{F}_{\mathrm{thermo}} = (G, C)$ is the thermodynamic fibre
encoding government intervention volume $G$ and institutional friction
$C$. The fibre acts as an \textbf{external field} that deforms the
metric parameters of the base manifold rather than adding coordinate
dimensions to it; this preserves general covariance of the base theory.

\subsection{Effective Liquidity Limit}

The single point of contact between the fibre and the base manifold is
a deformation of the liquidity absorption limit $c \to c_{\mathrm{eff}}$:
\begin{equation}
    \label{eq:ceff}
    \boxed{c_{\mathrm{eff}}(G,C) = c_0 \cdot
    \exp\!\left(-\lambda_G \frac{G}{G_0}
               - \lambda_C \frac{C}{C_0}\right)},
\end{equation}
where $c_0$ is the baseline limit calibrated on the pre-intervention
period, $G_0$ and $C_0$ are normalisation constants, and $\lambda_G$,
$\lambda_C$ are sensitivity parameters requiring empirical calibration.
The exponential form is motivated by the multiplicative suppression of
the liquid transaction space: each unit of intervention compounds the
deformation rather than adding to it linearly.

All results from Sections~3--9 carry over with $c \to c_{\mathrm{eff}}$;
in particular, the Lyapunov time becomes:
\begin{equation}
    \tau_L^{\mathrm{ext}} = \tau_0
    \sqrt{1 - \frac{v^2}{c_{\mathrm{eff}}(G,C)^2}}.
\end{equation}
The market reaches the Lyapunov singularity ($\tau_L^{\mathrm{ext}}
\to 0$) either when the market velocity $v \to c_{\mathrm{eff}}$ from
below (velocity-driven) or when $c_{\mathrm{eff}}$ is compressed down
to $|v|$ by intervention growth (field-driven). These are two
geometrically distinct paths to the same singularity condition
$v = c_{\mathrm{eff}}$.

\subsection{First Law: Energy Balance of the Extended Manifold}

The standard GISA conserved energy $E = A\gamma c$ must be augmented
with external flux terms:
\begin{equation}
    dE_{\mathcal{M}} = \delta Q_{\mathrm{ext}}
                     - \delta W_{\mathrm{market}}
                     + \delta\Phi_{\mathrm{corrupt}},
\end{equation}
where $\delta Q_{\mathrm{ext}}$ is the heat flux from the government
fiscal reservoir (subsidies, guaranteed purchases, rate suppression),
$\delta W_{\mathrm{market}}$ is the work extracted by real market
transactions, and $\delta\Phi_{\mathrm{corrupt}}$ is a dissipative
flux due to regulatory friction and institutional barriers. This term
is not conservative: it removes energy from the productive transaction
space without depositing it elsewhere in $\mathcal{M}$.

\subsection{Second Law: Entropy Export}

The market entropy $S_{\mathcal{M}}$ obeys the inequality from
Section~\ref{sec:entropy}. In the extended system the total entropy
budget splits:
\begin{equation}
    \frac{dS_{\mathrm{total}}}{dt}
    = \underbrace{\frac{dS_{\mathcal{M}}}{dt}}_{\leq\,0
      \text{ (bubble compression)}}
    + \underbrace{\frac{dS_{\mathrm{gov}}}{dt}}_{>\,0
      \text{ (sovereign debt, institutional fragility)}}
    \geq 0.
\end{equation}
Government intervention purchases a local reduction in market entropy
--- suppressing volatility and preventing price discovery --- at the
cost of accumulating entropy on the sovereign balance sheet.
The Second Law is not violated; its consequences are deferred to a
different reservoir and a later time interval.

\subsection{Lattice Defect Layer}

Equation~(\ref{eq:ceff}) describes a uniform field deformation. Real
institutional distortions are localised: a specific zoning regulation
blocks transactions in one price-volume neighbourhood; a developer
insolvency collapses one construction segment while the aggregate index
holds. These are \textbf{lattice defects} by analogy with condensed
matter physics.

The market Burgers vector measures the deviation of the real trajectory
from the geodesic over a closed loop $\mathcal{C}$ in
$(\ln V, \ln P)$ space:
\begin{equation}
    \mathbf{b}_{\mathrm{market}}
    = \oint_{\mathcal{C}} d(\ln P)
    - \oint_{\mathcal{C}} d(\ln P)_{\mathrm{geodesic}}.
\end{equation}
A nonzero $\mathbf{b}_{\mathrm{market}}$ is the geometric signature of
a dislocation in the manifold. Three defect classes and their market
realisations are:

\begin{center}
\begin{tabular}{p{3.2cm}p{3.2cm}p{6.0cm}}
\toprule
Defect type & Physics & Market realisation \\
\midrule
Point defect (vacancy) & Missing lattice atom &
  Blocked transaction (permit delay, zoning) \\
Line defect (dislocation) & Shifted atomic plane &
  Systemic corruption in one sector \\
Planar defect & Domain boundary &
  Two-speed market: subsidised vs.\ free \\
Interstitial (impurity) & Foreign atom &
  Institutional investor (REIT) with non-market logic \\
\bottomrule
\end{tabular}
\end{center}

Each defect contributes a local singularity to $c_{\mathrm{eff}}(x)$:
\begin{equation}
    c_{\mathrm{eff}}(x) = c_0
    \exp\!\left(-\lambda_G \frac{G}{G_0} - \lambda_C \frac{C}{C_0}\right)
    \cdot \left(1 - \sum_k
    \frac{\alpha_k}{\lvert\mathbf{x} - \mathbf{x}_k\rvert}\right),
\end{equation}
where $\mathbf{x}_k$ is the location of the $k$-th defect in
$(\ln V, \ln P)$ space and $\alpha_k$ its strength. Near a defect,
$c_{\mathrm{eff}} \to 0$ \emph{locally}, causing a segment-level
Lyapunov singularity without any global-index signal. This explains
the empirically observed pattern in which a pre-construction segment
collapses while the aggregate composite index remains stable.

\subsection{Stochastic Layer: Langevin Geodesic}

The deterministic geodesic equation is extended to a \textbf{covariant
Langevin equation}:
\begin{equation}
    \frac{D^2 x^\mu}{d\tau^2}
    = -\Gamma^\mu_{\mathrm{diss}}\frac{dx^\mu}{d\tau}
    + \eta^\mu(\tau),
\end{equation}
where $D^2/d\tau^2$ is the covariant second derivative (Christoffel
symbols included automatically), $\Gamma^\mu_{\mathrm{diss}}$ is a
dissipation tensor proportional to transaction costs, and
$\eta^\mu(\tau)$ has three statistically distinct components:
\begin{equation}
    \eta^\mu(\tau)
    = \eta^\mu_{\mathrm{micro}}(\tau)
    + \eta^\mu_{\mathrm{gov}}(\tau)
    + \eta^\mu_{\mathrm{defect}}(\tau).
\end{equation}

\begin{center}
\begin{tabular}{p{2.6cm}p{3.6cm}p{2.4cm}p{3.0cm}}
\toprule
Component & Source & Statistics & Existing link \\
\midrule
$\eta^\mu_{\mathrm{micro}}$ & Microstructure noise &
  Gaussian, $D_{\mathrm{micro}}$ & \S9.1 market vacuum \\
$\eta^\mu_{\mathrm{gov}}$ & Fiscal shock (program launch/cut) &
  L\'evy, $\alpha_L \in (1,2)$ &
  $\mathcal{F}_{\mathrm{thermo}}$, $\dot{G}$ \\
$\eta^\mu_{\mathrm{defect}}$ & Lattice collapse event &
  Poisson & Burgers vector $\mathbf{b}_{\mathrm{market}}$ \\
\bottomrule
\end{tabular}
\end{center}

The covariant noise correlation is:
\begin{equation}
    \langle \eta^\mu(\tau)\,\eta^\nu(\tau') \rangle
    = 2D\,g^{\mu\nu}(x)\,\delta(\tau - \tau'),
\end{equation}
with diffusion coefficient linked to the Unruh-analogue temperature
of Section~9.2:
\begin{equation}
    D = \frac{\kappa}{2\pi}\cdot\frac{\gamma^3}{A}
    \left\lvert\frac{dv}{d\tau}\right\rvert.
\end{equation}
The fluctuation-dissipation theorem then requires:
\begin{equation}
    \Gamma^\mu_{\mathrm{diss}}
    = \frac{D\,g^{\mu\mu}}{k_B T_{\mathcal{M}}},
\end{equation}
closing the thermodynamic consistency of the stochastic layer.
In real estate, transaction costs (land transfer tax, legal fees,
brokerage) reach 3--6\% of transaction value --- an order of magnitude
above the cryptocurrency fee tier of 0.35--0.75\% --- producing
substantially stronger geodesic damping.

\subsection{Observable Indicators and Monitoring Protocol}

The three-layer architecture produces three independent observable
signals, each sensitive to a different failure mode:

\begin{center}
\begin{tabular}{p{2.4cm}p{4.8cm}p{2.4cm}p{2.8cm}}
\toprule
Layer & Signal & Alert threshold & Data source \\
\midrule
1 (Thermo\-dynamic) &
  $c_{\mathrm{eff}}(t)$ monotone decrease &
  $c_{\mathrm{eff}} < 0.10$ &
  Federal budget housing expenditure \\
2 (Defect) &
  $\lvert\mathbf{b}_{\mathrm{market}}\rvert$ growth:
  MLS\,HPI vs.\ CMHC new-construction divergence &
  Sustained $> 5\%$ &
  CREA\,/\,CMHC monthly \\
3 (Stochastic) &
  Excess kurtosis of monthly $\Delta P$ &
  Kurtosis $> 6$ (L\'evy onset) &
  TRREB monthly \\
\bottomrule
\end{tabular}
\end{center}

All three signals are independent of the full calibration of
$\lambda_G$ and $\lambda_C$ and can be computed directly from
publicly available data. Synchronisation of all three signals
constitutes a multi-layer confirmation of proximity to the
Lyapunov singularity.

\subsection{Freeze Point: The $v = c_{\mathrm{eff}}$ Crossing}

The market freeze condition replaces the velocity-explosion singularity
of the flat-space model. Rather than $v$ accelerating toward $c_0$,
the freeze is reached when the declining $c_{\mathrm{eff}}(G, C)$
descends to meet the current market velocity $|v|$ from above. On the
$(G, C)$ parameter plane this defines a \textbf{freeze curve}:
\begin{equation}
    \boxed{v = c_0 \cdot \exp\!\left(
    -\lambda_G \frac{G}{G_0} - \lambda_C \frac{C}{C_0}\right)},
\end{equation}
whose location shifts with each fiscal cycle. Monitoring whether the
system approaches or recedes from this curve --- using Layer~1 and
Layer~3 signals as leading indicators --- converts the GISA engine
from a single-event detector into a continuous strategic surveillance
instrument for structural market fragility.

\subsection{Scope and Limitations of the Real-Market Extension}

\textbf{Important caveat.}\quad
This section constitutes a \emph{theoretical extension framework},
not a validated predictive model. The parameters
$\lambda_G$, $\lambda_C$, and $\alpha_k$ have not been empirically
calibrated against historical data. The framework is presented as:
\begin{enumerate}
    \item A research agenda for future empirical work (\S\ref{sec:future});
    \item A demonstration that the GISA geometric structure admits
          consistent thermodynamic extensions for open systems;
    \item A source of qualitative, directional analysis of structural
          fragility --- not quantitative point forecasts.
\end{enumerate}
All numerical estimates in \S12.3--\S12.9 are illustrative computations
under assumed parameter values. The directional conclusions (system in
`REPULS' regime; freeze condition approached via $c_{\mathrm{eff}}$
compression rather than velocity explosion; Layer~2 dislocation signal
empirically active in pre-construction segment) are robust to reasonable
parameter variation. The specific timelines and price targets are not.

\subsection{Empirical Calibration Programme}

The following quantities require numerical calibration before the
extended model can produce quantitative predictions:

\begin{enumerate}
    \item $c_0$: Golden-section fit to $(\ln V, \ln P)$ time series
          on the pre-intervention baseline (2003--2019 CREA data).
    \item $\lambda_G$, $\lambda_C$: Cross-sectional regression of
          observed $c_{\mathrm{eff}}$ (rolling golden-section fit)
          against $G/G_0$ and $C/C_0$ across multiple metropolitan
          areas and fiscal cycles.
    \item $\alpha_k$: Estimated from the magnitude of MLS HPI /
          CMHC new-construction price divergence in each identified
          segment.
    \item $\alpha_L$ (L\'evy index): Maximum-likelihood fit to the
          tail of the monthly $\Delta P$ distribution over the
          intervention period.
\end{enumerate}

Until this calibration is completed, all quantitative outputs from
the extended model (freeze dates, threshold prices) are to be treated
as qualitative structural analysis rather than precision forecasts.

\appendix

\section{C++ Implementation Overview}

Production-ready, header-only, C++17, no external dependencies:

\begin{itemize}
    \item \texttt{metric.hpp} --- \texttt{LorentzMetric}: $ds^2$,
          $\gamma(v,c)$, geodesic price/log-price helpers,
          \texttt{ln\_price\_from\_geodesic}.
    \item \texttt{curvature.hpp} --- \texttt{RicciScalar}: non-uniform
          central differences, template dispatch for 9- and 10-tuple
          history, $c$-acceleration fallback for sparse data.
    \item \texttt{manifold.hpp} --- \texttt{ManifoldTracker}: per-bar
          state update; $c$-stability convergence; $c$-collapse detector
          (prev\_c captured before last\_c\_ is overwritten); curvature-
          before-deviation classification order; adaptive thresholds for
          $R$ and $\xi$ that suppress early false positives; ds2
          sign-change breaking check.
    \item \texttt{geodesic.hpp} --- \texttt{GeodesicTracker}: log-space
          $A$-drift (\texttt{|lnP - entry\_lnA|}); \texttt{HOLD}/
          \texttt{PREPARE\_EXIT}/\texttt{SELL}; \texttt{set\_entry}
          seeded from first bar's $c$ to avoid false collapse detection
          at bar 0.
    \item \texttt{gisa.hpp} --- master include.
\end{itemize}

\section*{Keywords}

Exponential impulse; Lorentz factor; Lorentzian manifold; market
microstructure; VWAP; metric tensor; Christoffel symbols; Ricci scalar;
geodesic deviation; golden ratio; Nash equilibrium; Lyapunov time; entropy
bound; dimensionality collapse; holographic principle; AdS/CFT;
Bekenstein entropy; Unruh temperature; Hawking radiation; effective-regime
theory; dynamic $c$-calibration; golden-section search; Savitzky--Golay;
out-of-sample validation.

\section*{Acknowledgments}

Validated on real exchange data: LTC-USDC (May~22--23, 2026), FET-USDC
(March~19--20, 2026), BICO-USDC (June~19--20, 2026).

\end{document}